% Unit 120 terjemahan Bahasa Indonesia dari chapter8.tex baris 2108--2214.
% Sumber: Methods of Algebra, Volume 2, commit
% 9a5803ff77dd3257484cb177f851a73770a59dd3 (CC BY 4.0).
% stable-unit-id: o014.aljabr2.chapter8.exercises-b
% Cakupan: kelanjutan latihan tentang objek L-projektif hingga akhir bab.
% Koreksi sumber terpaksa: baris 2113 identity_{X_n} diperjelas sebagai
% identity_{L^{n+1} M}; baris 2171 objek homologi berada dalam A, bukan Ch(A).

% segment-id: o014.li2.chapter8.exercises-b.p001
	Dengan dualitas, pembahasan berikut terutama berfokus pada konsep
	$L$-projektif. Dengan metode Contoh~\ref{eg:monad-simplicial}, dalam
	$\End(\mathcal{D})$ diperoleh objek simplisial teraugmentasi
	$(L^{n+1})_{n\geq-1}$; data $d_i$, $s_j$, dan sebagainya dihilangkan dari
	notasi. Untuk objek $L$-projektif $\mathcal{M}$ dan $f$ yang
	bersesuaian, definisikan
	% segment-id: o014.li2.chapter8.exercises-b.q001
	\[ k_n := L^{n+1} f : L^{n+1} \mathcal{M} \to L^{n+2} \mathcal{M}, \quad n \geq -1. \]
	Buktikan bahwa ini membuat objek simplisial teraugmentasi
	$(L^{n+1}\mathcal{M})_{n\geq-1}$ dapat dikontraksi dari kanan
	(Definisi~\ref{def:contractible-aug}).

	% segment-id: o014.li2.chapter8.exercises-b.hint001
	\begin{hint}
		Kita tahu bahwa
		$\epsilon_{\mathcal{M}}k_{-1}=\epsilon_{\mathcal{M}}f
		=\identity_{\mathcal{M}}$. Menerapkan $L^{n+1}$ pada kedua ruas
		memberikan
		$d_{n+1}k_n=\identity_{L^{n+1}\mathcal{M}}$, sedangkan naturalitas
		memastikan bahwa
		% segment-id: o014.li2.chapter8.exercises-b.q002
		\[\begin{tikzcd}
			L\mathcal{M} \arrow[r, "Lf"] \arrow[d, "{\epsilon_{\mathcal{M}}}"'] & L^2 \mathcal{M} \arrow[d, "{\epsilon_{L\mathcal{M}}}"] \\
			\mathcal{M} \arrow[r, "f"'] & L\mathcal{M}
		\end{tikzcd} \quad\text{komutatif, yakni}\; d_0 k_0 = f\epsilon_{\mathcal{M}} = k_{-1} \epsilon_{\mathcal{M}}. \]
		Ganti $f$ dengan $L^{n-i}f$, lalu terapkan $L^i$ pada diagram yang
		bersesuaian. Dengan cara yang sama diperoleh
		$d_ik_n=k_{n-1}d_i$ untuk $0\leq i\leq n$.
	\end{hint}

	% segment-id: o014.li2.chapter8.exercises-b.ex001
	\item Untuk suatu gelanggang $R$, pasangan adjoin bebas--pelupa
	$\cate{Set}\leftrightarrows R\dcate{Mod}$ menentukan komonad $L$ pada
	$R\dcate{Mod}$. Buktikan bahwa modul-$R$ kiri bersifat $L$-projektif
	jika dan hanya jika ia merupakan modul projektif.

	% segment-id: o014.li2.chapter8.exercises-b.ex002
	\item Tinjau pasangan adjoin
	% segment-id: o014.li2.chapter8.exercises-b.q003
	$\begin{tikzcd}
		F: \mathcal{C} \arrow[shift left, r] & \mathcal{D}: G \arrow[shift left, l]
	\end{tikzcd}$
	dan komonad yang bersesuaian
	$(L,\delta,\epsilon)=(FG,F\eta G,\epsilon)$ pada
	$\End(\mathcal{D})$ (Contoh~\ref{eg:adjunction-monad}). Buktikan bahwa
	objek berbentuk $F\mathcal{N}$, untuk
	$\mathcal{N}\in\Obj(\mathcal{C})$, selalu $L$-projektif. Sebagai
	akibatnya, untuk setiap $\mathcal{M}\in\Obj(\mathcal{D})$, setiap suku
	dari objek simplisial $(L^{n+1}\mathcal{M})_{n\geq0}$ merupakan objek
	$L$-projektif.

	% segment-id: o014.li2.chapter8.exercises-b.hint002
	\begin{hint}
		Ambil
		$f=F\eta_{\mathcal{N}}:F(\mathcal{N})\to FGF(\mathcal{N})
		=L(F(\mathcal{N}))$. Maka
		$\epsilon_{F(\mathcal{N})}f=\identity_{F(\mathcal{N})}$ tepat
		merupakan salah satu identitas segitiga pasangan adjoin.
	\end{hint}

	% segment-id: o014.li2.chapter8.exercises-b.ex003
	\item Tetap tinjau komonad $L$ yang berasal dari pasangan adjoin
	$(F,G)$. Tunjukkan bahwa $\mathcal{P}\in\Obj(\mathcal{D})$ bersifat
	$L$-projektif jika dan hanya jika ia memiliki sifat pengangkatan
	berikut. Untuk setiap morfisme
	$\alpha:\mathcal{M}_1\to\mathcal{M}_2$ dan
	$\phi:\mathcal{P}\to\mathcal{M}_2$ dalam $\mathcal{D}$, jika
	$G\alpha$ memiliki invers kanan, maka terdapat
	$\beta:\mathcal{P}\to\mathcal{M}_1$ sedemikian rupa sehingga
	$\phi=\alpha\beta$.

	% segment-id: o014.li2.chapter8.exercises-b.hint003
	\begin{hint}
		Untuk arah ``jika'', terapkan sifat pengangkatan pada
		$\alpha:=\epsilon_{\mathcal{P}}:FG(\mathcal{P})\to\mathcal{P}$ dan
		$\phi:=\identity$. Untuk arah ``hanya jika'', jelaskan terlebih dahulu
		bahwa $\mathcal{P}$ dalam sifat pengangkatan dapat diganti oleh
		$FG(\mathcal{P})$, lalu gunakan sifat adjoin untuk menunjukkan bahwa
		% segment-id: o014.li2.chapter8.exercises-b.q004
		\[ \alpha_*:\Hom(FG(\mathcal{P}),\mathcal{M}_1)\to
		\Hom(FG(\mathcal{P}),\mathcal{M}_2) \]
		surjektif.
	\end{hint}

	% segment-id: o014.li2.chapter8.exercises-b.ex004
	\item Untuk gelanggang komutatif $\Bbbk$ sebarang dan homomorfisme
	aljabar-$\Bbbk$ $S\to R$, tinjau pasangan-pasangan adjoin
	% segment-id: o014.li2.chapter8.exercises-b.q005
	\begin{equation*}\begin{tikzcd}[row sep=small]
		R \dotimes{S} (\cdot): S\dcate{Mod} \arrow[shift left, r] & R\dcate{Mod}: \text{lupa}, \arrow[shift left, l] \\
		\text{lupa}: R\dcate{Mod} \arrow[shift left, r] & S\dcate{Mod}: \Hom_S(R, \cdot) \arrow[shift left, l]
	\end{tikzcd}\end{equation*}
	yang masing-masing memberikan komonad $L$ dan monad $T$ pada
	$R\dcate{Mod}$. Gunakan sifat pengangkatan di atas, atau dualnya, untuk
	tentukan syarat perlu dan cukup bagi modul-$R$ kiri $M$ untuk menjadi
	objek $L$-projektif (atau $T$-injektif), lalu bandingkan dengan modul
	projektif (atau injektif).

	Atas dasar ini dapat dibangun teori yang disebut \emph{aljabar homologis
	relatif}; lihat \cite{Ho56}. Ciri khasnya ialah bahwa teori tersebut
	hanya meninjau barisan eksak modul-$R$ yang eksak terbelah di atas $S$.
	Latihan berikut akan menyelidiki versi relatif funktor $\Tor$ dan $\Ext$.
	\index{aljabar homologis relatif@aljabar homologis relatif (relative homological algebra)}

	% segment-id: o014.li2.chapter8.exercises-b.ex005
	\item Tinjau pasangan adjoin di antara kategori abelian
	% segment-id: o014.li2.chapter8.exercises-b.q006
	$\begin{tikzcd}
		F: \mathcal{C} \arrow[shift left, r] & \mathcal{D}: G \arrow[shift left, l]
	\end{tikzcd}$.
	Buktikan untuk setiap $\mathcal{M}\in\Obj(\mathcal{D})$ bahwa:
	\begin{enumerate}[(i)]
		\item Terdapat objek $\mathcal{P}_\bullet$ dari
		$\cate{Ch}_{\geq0}(\mathcal{D})$ beserta morfisme
		$\epsilon:\mathcal{P}_\bullet\to\mathcal{M}$ sedemikian rupa sehingga
		\begin{compactitem}
			\item setiap $\mathcal{P}_n$ bersifat $L$-projektif untuk $n\geq0$;
			\item morfisme identitas pada
			$\cdots\to G\mathcal{P}_1\to G\mathcal{P}_0\to
			G\mathcal{M}\to0$ homotopik nol; dengan kata lain, ini kompleks
			eksak terbelah.
		\end{compactitem}
		\begin{hint}
			Terapkan Proposisi~\ref{prop:bar-null-homotopic}.
		\end{hint}
		\item Untuk setiap dua morfisme
		$\epsilon:\mathcal{P}_\bullet\to\mathcal{M}$ dan
		$\epsilon':\mathcal{P}'_\bullet\to\mathcal{M}$ yang memenuhi
		syarat-syarat di atas, terdapat ekuivalensi homotopi
		$f:\mathcal{P}\to\mathcal{P}'$ sedemikian rupa sehingga
		$\epsilon'f=\epsilon$.
		\begin{hint}
			Gunakan sifat pengangkatan yang diperkenalkan di atas.
		\end{hint}
	\end{enumerate}

	\index{homologi komonad@homologi komonad (cotriple homology)}
	\index{konstruksi bar@konstruksi bar}
	% segment-id: o014.li2.chapter8.exercises-b.ex006
	\item (M.\ Barr, J.\ M.\ Beck) Misalkan $L$ sebuah komonad pada
	kategori $\mathcal{D}$. Tuliskan $\mathrm{Bar}(\mathcal{M})$ untuk
	objek simplisial yang diperoleh dengan mengevaluasi konstruksi
	Contoh~\ref{eg:monad-simplicial} pada
	$\mathcal{M}\in\Obj(\mathcal{D})$. Untuk setiap kategori abelian
	$\mathcal{A}$ dan funktor $E:\mathcal{D}\to\mathcal{A}$, definisikan
	\emph{homologi komonad} dari $(L,E)$ sebagai keluarga objek dalam
	$\mathcal{A}$
	% segment-id: o014.li2.chapter8.exercises-b.q007
	\[ \Hm^L_n(E; \mathcal{M}) := \Hm_n\left( \mathrm{C}(E\mathrm{Bar}(\mathcal{M})) \right), \quad n \in \Z_{\geq 0}, \]
	dengan $E\mathrm{Bar}(\mathcal{M})$ berarti menerapkan $E$ pada setiap
	suku $\mathrm{Bar}(\mathcal{M})$. Secara dual, jika $T$ monad pada
	kategori $\mathcal{C}$ dan $F:\mathcal{C}\to\mathcal{A}$ funktor menuju
	kategori abelian, kita dapat mendefinisikan secara serupa nilai
	\emph{kohomologi monad} $(T,F)$ pada
	$\mathcal{N}\in\Obj(\mathcal{C})$; perinciannya tidak kita uraikan.

	\begin{enumerate}[(i)]
		\item Tuliskan morfisme alami
		$\Hm^L_0(E;\mathcal{M})\to E\mathcal{M}$ dan tunjukkan bahwa ia
		merupakan isomorfisme jika $\mathcal{M}$ objek $L$-projektif.
		\item Misalkan $L$ berasal dari pasangan adjoin $(F,G)$ di antara
		kategori abelian dan $E$ funktor aditif. Buktikan bahwa jika
		$\mathcal{P}_\bullet\to\mathcal{M}$ morfisme dalam
		$\cate{Ch}_{\geq0}(\mathcal{D})$ yang memenuhi
		\begin{compactitem}
			\item setiap $\mathcal{P}_n$ bersifat $L$-projektif untuk $n\geq0$;
			\item morfisme identitas pada
			$\cdots\to G\mathcal{P}_1\to G\mathcal{P}_0\to
			G\mathcal{M}\to0$ homotopik nol,
		\end{compactitem}
		maka
		$\Hm^L_n(E;\mathcal{M})\simeq\Hm_n(E\mathcal{P}_\bullet)$.
		Jadi homologi komonad dapat dihitung menggunakan resolusi
		$L$-projektif semacam itu.
		\item Lanjutkan dengan asumsi di atas. Sebuah pasangan morfisme
		$\mathcal{M}'\to\mathcal{M}\to\mathcal{M}''$ dalam $\mathcal{D}$
		disebut barisan eksak pendek \emph{$G$-terbelah} jika setelah
		menerapkan $G$ ia menjadi barisan eksak pendek terbelah dalam
		$\mathcal{C}$. Buktikan bahwa barisan eksak pendek $G$-terbelah
		secara alami menginduksi barisan eksak panjang untuk
		$\Hm^L_n(E;\mathord\cdot)$.
	\end{enumerate}

	\index{funktor Ext!relatif (relative)}
	\index{funktor Tor!relatif (relative)}
	% segment-id: o014.li2.chapter8.exercises-b.ex007
	\item (G.\ Hochschild \cite{Ho56}) Untuk sebarang gelanggang komutatif
	$\Bbbk$ dan homomorfisme aljabar-$\Bbbk$ $S\to R$, tinjau komonad $L$
	pada $R\dcate{Mod}$ yang ditentukan oleh pasangan adjoin
	% segment-id: o014.li2.chapter8.exercises-b.q008
	\[\begin{tikzcd}
		R \dotimes{S} (\cdot): S\dcate{Mod} \arrow[shift left, r] & R\dcate{Mod}: \text{lupa} \arrow[shift left, l]
	\end{tikzcd}\]
	Untuk modul-$R$ kanan $A$ dan modul-$R$ kiri $B,N$, definisikan
	modul-$\Bbbk$
	% segment-id: o014.li2.chapter8.exercises-b.q009
	\begin{align*}
		\Tor^{R|S}_n(A, N) & := \Hm^L_n\left( A \dotimes{R}(\cdot) ; N \right), \\
		\Ext_{R|S}^n(N, B) & := \Hm^L_n\left( \Hom_R(\cdot, B) ; N \right),
	\end{align*}
	dengan $\Hom_R(\mathord\cdot,B)$ dipandang sebagai funktor
	$R\dcate{Mod}\to S\dcate{Mod}^{\opp}$. Keduanya disebut funktor
	$\Tor$ relatif dan $\Ext$ relatif.
	\begin{enumerate}[(i)]
		\item Tunjukkan bahwa keduanya funktorial dalam setiap variabel.
		Uraikan secara eksplisit kompleks rantai yang bersesuaian, dan
		tunjukkan bahwa
		% segment-id: o014.li2.chapter8.exercises-b.q010
		\[ \Tor^{R|S}_0(A, N) \simeq A \dotimes{R} N, \quad \Ext_{R|S}^0(N, B) \simeq \Hom_R(N, B). \]

		\item Tunjukkan bahwa bifunktor $\Tor^{R|S}_n$ memiliki sifat
		``keseimbangan'' yang serupa dengan
		Teorema~\ref{prop:balanced-primer}.\footnote{$\Ext_{R|S}^n$ tidak
		dibahas karena melibatkan konsep kohomologi monad, yang tidak kita
		uraikan di sini; lihat rujukan untuk perinciannya.}

		\item Buktikan bahwa jika $I$ ideal dua sisi dari $S$ dan $R=S/I$,
		maka untuk $n>0$,
		$\Tor_n^{R|S}=0=\Ext^n_{R|S}$. Jadi $\Tor$ relatif dan $\Ext$
		relatif berbeda dari versi absolut dalam \S\ref{sec:Ext-Tor}.
		\begin{hint}
			Dalam hal ini $L=\identity_{R\dcate{Mod}}$.
		\end{hint}

		\item Tuliskan $R^e=R\dotimes{\Bbbk}R^{\opp}$. Buktikan bahwa
		$\HHm_n(M)\simeq\Tor^{R^e|\Bbbk}_n(M,R)$ dan
		$\HHm^n(M)\simeq\Ext_{R^e|\Bbbk}^n(R,M)$; lihat
		Contoh~\ref{eg:HH-comonad}.
		\begin{hint}
			Reduksikan persoalan ke pernyataan bahwa
			$\mathrm{C}(\mathrm{Bar}(R))_n=R\otimes R^{\otimes n}\otimes R$
			merupakan modul-$R^e$ yang $L$-projektif.
		\end{hint}
	\end{enumerate}
\end{Exercises}
